\section{Discusión}
Para Venezuela, el punto de entrada puede ser una pregunta local que conecte una asignatura, una comunidad y un conjunto de recursos abiertos. La escala inicial debe corresponder al tiempo de tutoría y revisión disponible. Materiales descargables y tareas asincrónicas permiten acomodar conectividad heterogénea; encuentros breves sirven para resolver dificultades y comparar decisiones. La coordinación requiere compromisos concretos de personas e instituciones, además de visibilidad pública.

La comparación muestra que preservar un producto y sostener una comunidad exigen decisiones diferentes. Un historial permite recuperar una edición, pero la revisión necesita personas; una tesis reconoce un aprendizaje, pero su reutilización necesita acceso al procedimiento. Un hackatón puede abrir una trayectoria si ofrece nuevas responsabilidades y conexiones. En consecuencia, la continuidad debe incorporarse al presupuesto y al reparto de tareas del encuentro.

\paragraph{Alcance y límites} La selección de casos es intencional y cercana a la trayectoria de los autores. Las fuentes describen productos con unidades y períodos distintos; faltan registros comparables por edición y seguimiento longitudinal. La matriz es una síntesis analítica de mecanismos, no una estimación causal del aprendizaje ni una clasificación de eficacia. La equivalencia de cohortes entre las noticias de 2023 permanece sin confirmar. No clasificamos la definición de un formato como más estricta que la de otro.
